The Kerr-Schild coordinates are labels of the flat background and $H$, $k$ those of the Kerr gauge field $\hb(b) = b + H(b \cdot k)k$, whose derived metric is $\eta - 2H\,k \otimes k$ (Fact \claimref{synthesis.imported-gtg-dictionary}(i)). \emph{Provenance.} The identity $z^2 = w$ and the form $H = M\operatorname{Re}(1/z)$ are the classical representation of the Kerr solution as the field of a point source displaced to an imaginary position (Appell, Math.\ Ann.\ 30 (1887); Newman and Janis, J.\ Math.\ Phys.\ 6 (1965) 915; Israel, Phys.\ Rev.\ D 2 (1970) 641; Newman, J.\ Math.\ Phys.\ 14 (1973) 102; cited from memory), and the inverse-cube scale is the classical $\Psi_2$ of Fact \claimref{interior.imported-kerr-curvature}. They are verified below in the store's conventions; the content added is clauses (3) and (4) read on the double cover.

\emph{(1)} With $x^2 + y^2 = (r^2 + a^2)\sin^2\theta$ and $Z = r\cos\theta$, $w = (r^2 + a^2)\sin^2\theta + r^2\cos^2\theta - 2Iar\cos\theta - a^2 = r^2 - a^2\cos^2\theta - 2Iar\cos\theta = (r - Ia\cos\theta)^2$, using $I^2 = -1$ in the commutative subalgebra $\R \oplus \R I \cong \C$ of Proposition \claimref{interior.complex-factorisation-location-ring}(1). The real and pseudoscalar parts of $w = \varrho^2 + Z^2 - a^2 - 2IaZ$ vanish together exactly when $Z = 0$ and $\varrho = a$; there $\partial_\varrho w = 2\varrho = 2a$ and $\partial_Z w = 2Z - 2Ia = -2Ia$, real and pseudoscalar respectively and both nonzero, so $w$ is a submersion onto $\C$ near the ring and a complex normal coordinate. Hence $z$ with $z^2 = w$ is the coordinate of the double cover branched along the ring, and $z = 0$ is the ring, matching Proposition \claimref{interior.complex-factorisation-location-ring}(3). The involution $(r, \theta) \mapsto (-r, \pi - \theta)$ leaves $(r^2 + a^2)\sin^2\theta$ and $r\cos\theta$ unchanged and sends $z = r - Ia\cos\theta$ to $-r + Ia\cos\theta = -z$.

\emph{(2)} By Fact \claimref{interior.imported-kerr-schild-form-kerr-gauge}, $H = Mr^3/(r^4 + a^2\cos^2\theta\,r^2) = Mr/(r^2 + a^2\cos^2\theta)$ for $r \neq 0$, and $1/z = \tilde z/(z\tilde z) = (r + Ia\cos\theta)/(r^2 + a^2\cos^2\theta)$, whose real part is $r/(r^2 + a^2\cos^2\theta)$; on the disc $r = 0$, $\cos\theta \neq 0$, both sides vanish. So $H = M\operatorname{Re}(1/z) = M\cos(\arg z)/|z|$, zero exactly when $z$ is purely pseudoscalar, i.e. on $r = 0$, whose two faces ($\cos\theta > 0$ and $< 0$) are $\arg z = -\pi/2$ and $+\pi/2$. Downstairs $|z| = |w|^{1/2}$ and $\arg z = \tfrac12\arg w$ on a branch. The metric is $g = \eta - 2Hk\otimes k$ by Proposition \claimref{interior.unipotency-across-kerr-family}(2). With the Kerr-Schild null vector of Fact \claimref{interior.imported-kerr-schild-form-kerr-gauge} in its standard components $\bigl(1, (rx + ay)/(r^2 + a^2), (ry - ax)/(r^2 + a^2), Z/r\bigr)$, $Z/r = \cos\theta$, the limit at the ring ($r \to 0$, $\cos\theta \to 0$) is $(1, y/a, -x/a, 0)$, the time direction plus the unit tangent of the circle $x^2 + y^2 = a^2$; the components are cited in their classical form \cite{KerrSchild1965}, only this limit is used, and the orientation of the limiting tangent depends on the sign convention for the $a$-terms of $k$.

\emph{(3)} $\operatorname{Re}(1/(-z)) = -\operatorname{Re}(1/z)$ and $(-z)^{-3} = -z^{-3}$, the latter being the scale of the curvature operator of Fact \claimref{interior.imported-kerr-curvature}, which therefore changes sign under the deck transformation.

\emph{(4)} On a loop around $z = 0$, $z^{-3}$ has winding number $-3$ (argument principle, or the exact crossing count of the check); in the identification of Proposition \claimref{interior.complex-scale-inverse-cube-killing-spinor} the phase of $\mathfrak M$ is the duality angle of the curvature eigenvalues (Remark \claimref{note.interior.duality-winding-around-ring}), so the duality phase turns by $-6\pi$. A loop once around the ring downstairs (closed in Kerr-Schild coordinate space, an open path between the sheets in the maximal extension) lifts to a path on the double cover along which $\arg z$ increases by $\pi$, from $z_0$ to $-z_0$, so the phase of $\mathfrak M$ turns by $-3\pi$ and $\mathfrak M(-z_0) = -\mathfrak M(z_0)$, which is the parity of (3): the holonomy $e^{-3\pi I} = -1$ is the deck transformation's action on the scale. The check confirms that the image of the double-cover loop under $z \mapsto z^2$ winds twice around $w = 0$.
